\section{Problem 788}

\subsection*{1) ROUND-2 OBJECTIVE}
Path \textbf{(A)}: prove a new unconditional upper bound. This is the most promising direction because Round 1 already reduced the problem to an independence-number problem in a graph, and there is now a strong external theorem on random Cayley sum graphs that can be transferred into this setting.

\subsection*{2) Round-1 FOUNDATION USED}
I use the following Round-1 results as fixed:
\begin{enumerate}
\item the graph reformulation
\[
 f(n)=\min_{B \subset (2n,4n)} \bigl(|B|+\alpha(G_B)\bigr),
\]
where $G_B$ has vertex set $I_n:=\{n+1,\dots,2n-1\}$ and an edge $xy$ whenever $x\neq y$ and $x+y \in B$;
\item the lower bound
\[
 f(n) \ge 2\sqrt{n-1}-1;
\]
\item the small exact computations for tiny $n$, used only as a check.
\end{enumerate}

\subsection*{3) NEW INSIGHT / TOOL (ROUND-2)}
The new tool is a transference principle from random Cayley sum graphs to Problem 788.

The external input is Alon--Pham (2025): for an abelian group $G$ of size $N$ and edge-density parameter $p \le 1/2$, the random Cayley sum graph on $G$ has independence number
\[
 \widetilde O(p^{-3/2})
\]
with high probability.

I will prove a general reduction showing that any bound of the form $\widetilde O(p^{-c})$ for random Cayley sum graphs implies
\[
 f(n) \le n^{c/(c+1)+o(1)}.
\]
Substituting $c=3/2$ yields the new exponent $3/5$.

\subsection*{4) ATTACK PLAN (ROUND-2)}
The exact Round-1 gap was on the upper-bound side: we had only the older $2/3$ exponent. The plan is:
\begin{enumerate}
\item embed the interval graph $G_B$ into a Cayley sum graph on a cyclic group of prime order $q>4n$ so that there is no modular wrap-around;
\item choose the Cayley generating set randomly with density $p$;
\item use Alon--Pham to make the independence number small and Chernoff to keep the generating set sparse;
\item optimize $pn + p^{-c}$.
\end{enumerate}

\subsection*{5) WORK (ROUND-2)}
Recall from Round 1 that $f(n)=\min_B (|B|+\alpha(G_B))$.

\paragraph{Lemma 788.1 (general transfer principle).}
Assume the following external hypothesis for some fixed $c>0$:

\medskip
\noindent
for every finite abelian group $G$ of order $N$ and every $p \le 1/2$, the random Cayley sum graph on $G$ has independence number at most
\[
 \widetilde O(p^{-c})
\]
with high probability.
\medskip

Then
\[
 f(n) \le n^{c/(c+1)+o(1)}.
\]

\paragraph{Proof.}
Fix $n$ large. By Bertrand's postulate there is a prime $q$ with
\[
 4n < q < 8n.
\]
Work in the cyclic group $G = \mathbf Z / q\mathbf Z$.

Choose a random subset $S \subseteq G$ by including each element independently with probability $p$, where $p \le 1/2$ will be chosen later. Let $\Gamma_S^+$ be the associated random Cayley sum graph on $G$.

By the external hypothesis, with high probability,
\[
 \alpha(\Gamma_S^+) \le K p^{-c} (\log q)^K
\]
for some absolute constant $K$.
Also, by a standard Chernoff bound, with high probability,
\[
 |S| \le 2pq.
\]
Hence for large $q$ there exists at least one realisation of $S$ for which both estimates hold simultaneously.

Fix such an $S$, and define
\[
 B := S \cap \{2n+1,2n+2,\dots,4n-1\}.
\]
Then $B \subset (2n,4n)$, as required by the problem.

Now consider the interval of vertices
\[
 I_n = \{n+1,n+2,\dots,2n-1\} \subset G.
\]
If $x,y \in I_n$ are distinct, then
\[
 2n+3 \le x+y \le 4n-3 < q,
\]
so the sum $x+y$ in $G$ is the same as the ordinary integer sum. Therefore, for distinct $x,y \in I_n$,
\[
 x+y \in S \iff x+y \in B.
\]
This means exactly that the graph $G_B$ from Problem 788 is the induced subgraph of $\Gamma_S^+$ on the vertex set $I_n$.

Every independent set in $G_B$ is therefore an independent set in $\Gamma_S^+$, and hence
\[
 \alpha(G_B) \le \alpha(\Gamma_S^+) \le K p^{-c} (\log q)^K.
\]
Also,
\[
 |B| \le |S| \le 2pq.
\]
Consequently,
\[
 f(n) \le |B| + \alpha(G_B)
 \le 2pq + K p^{-c} (\log q)^K.
\]
Since $q \asymp n$, this is
\[
 f(n) \le \widetilde O(pn + p^{-c}).
\]
Choose
\[
 p = n^{-1/(c+1)}.
\]
Then
\[
 pn = n^{c/(c+1)}
\qquad\text{and}\qquad
p^{-c} = n^{c/(c+1)}.
\]
Thus
\[
 f(n) \le \widetilde O\!\left(n^{c/(c+1)}\right)
 = n^{c/(c+1)+o(1)}.
\]
This proves the lemma.
\hfill$\square$

\paragraph{External theorem (Alon--Pham, 2025, Theorem 4).}
For an abelian group of size $N$ and $p \le 1/2$, the random Cayley sum graph has independence number at most
\[
 \widetilde O(p^{-3/2})
\]
with high probability.

\paragraph{Corollary 788.2.}
One has
\[
 f(n) \le n^{3/5+o(1)}.
\]

\paragraph{Proof.}
Apply Lemma 788.1 with $c=3/2$. Then
\[
 \frac{c}{c+1} = \frac{3/2}{5/2} = \frac35.
\]
So
\[
 f(n) \le n^{3/5+o(1)}.
\]
\hfill$\square$

\paragraph{Summary.}
Combining Corollary 788.2 with the Round-1 lower bound gives
\[
 n^{1/2} \ll f(n) \le n^{3/5+o(1)}.
\]
This improves the previously quoted upper exponent $2/3$.

\subsection*{6) ADVERSARIAL VERIFICATION}
\begin{enumerate}
\item The no-wrap condition is the key technical point, and it is valid because $q>4n$ while all relevant sums of distinct elements of $I_n$ lie in $[2n+3,4n-3]$.
\item I used only distinct pairs in the interval graph $G_B$, exactly as in the original problem. Under the standard simple-graph convention for Cayley sum graphs, the induced-subgraph identification is exact on distinct vertices.
\item The probabilistic argument is converted to a deterministic existence statement by intersecting two high-probability events: one for $\alpha(\Gamma_S^+)$ and one for $|S|$.
\item The optimisation is correct:
\[
 pn + p^{-c}
\]
is balanced by $p=n^{-1/(c+1)}$, giving exponent $c/(c+1)$.
\item Since $p=n^{-1/(c+1)} \to 0$, the assumption $p \le 1/2$ holds for all sufficiently large $n$; the finitely many small $n$ are absorbed into the implied constants.
\end{enumerate}

\subsection*{7) FINAL}
\textbf{UNRESOLVED (BUT STRICTLY ADVANCED).}

The strict advance is a rigorous new upper bound:
\[
 f(n) \le n^{3/5+o(1)}.
\]
This improves the older $2/3$ exponent and comes from a general transfer lemma from random Cayley sum graph independence bounds to Problem 788.

\subsection*{8) COMPLETION ESTIMATE}
\textbf{COMPLETION: 65\%}

\subsection*{9) REFERENCES}
\begin{thebibliography}{99}
\bibitem{Choi788}
S. L. G. Choi,
\emph{On a combinatorial problem in number theory},
Proc. London Math. Soc. (3) 23 (1971), 629--642.

\bibitem{BSS788}
A. Baltz, T. Schoen, and A. Srivastav,
\emph{Probabilistic construction of small strongly sum-free sets via large Sidon sets},
Colloq. Math. 86 (2000), 171--176.

\bibitem{AlonPham788}
N. Alon and H. T. Pham,
\emph{Random Cayley graphs and random sumsets},
arXiv:2509.02561 (2025).

\bibitem{Bloom788}
T. F. Bloom,
\emph{Erd\H{o}s Problem \#788}, updated January 2026,
\texttt{https://www.erdosproblems.com/788}.
\end{thebibliography}
